\section{Simulação do Movimento de uma Partícula}
\label{sec:simulacao}

Esta seção descreve a arquitetura computacional utilizada para simular o movimento da partícula sob ação da gravidade e do arrasto atmosférico. A ideia central do projeto é que \textbf{o estado completo do sistema em qualquer instante é determinado pelo triplo} $(m,\,\vec{r},\,\vec{v})$, e todas as forças são funções puras desse estado - sem memória e sem dependência explícita do tempo. Essa escolha de design torna o código modular: novas forças podem ser adicionadas à simulação sem modificar o laço de integração.

Durante essa sessão será exposto partes do código, mas ele esta aberto ao publico no repositório hospedado no github \cite{mecanica-geral-repo-2026}.

\subsection{Estado do Sistema e a Partícula}

O estado do sistema é representado pela partícula $P$, definida pelos seus atributos:
\begin{equation}
    P = (m,\,\vec{r},\,\vec{v}), \qquad
    \vec{r} = \begin{bmatrix}x\\y\\z\end{bmatrix}, \qquad
    \vec{v} = \begin{bmatrix}v_x\\v_y\\v_z\end{bmatrix}
\end{equation}

onde $m \in \mathbb{R}^+$ é a massa, $\vec{r} \in \mathbb{R}^3$ a posição e $\vec{v} \in \mathbb{R}^3$ a velocidade. A força total sobre a partícula é a soma de contribuições individuais, cada uma dependendo apenas de $P$:
\begin{equation}
    \vec{F}_\text{total}(P) = \sum_n \vec{F}_n(P)
\end{equation}

Cada $\vec{F}_n$ é uma função do tipo \texttt{ForceFunc}: recebe o estado atual $P$ e retorna um vetor em $\mathbb{R}^3$. A construção das forças como funções de $P$ garante que elas não dependem de variáveis globais mutáveis nem do instante $t$ explicitamente.

A classe \texttt{Particle} encapsula o estado e dois métodos de atualização usados na integração. O método \texttt{update\_velocity} calcula a resultante $\vec{F}_\text{total}$, obtém $\vec{a} = \vec{F}_\text{total}/m$ e aplica o passo de Euler à velocidade. O método \texttt{update\_position} aplica o passo de Euler à posição usando a velocidade \emph{já atualizada} - esse detalhe caracteriza o Euler simplético e será retomado na Sec.~\ref{subsec:integracao}.

\begin{lstlisting}[caption={Estado do sistema: tipo ForceFunc e classe Particle}, label={lst:particle}]
type ForceFunc = Callable[["Particle"], Vec3]

class Particle:
    mass: float
    position: Vec3
    velocity: Vec3
    def __init__(self,
        mass: float,
        position: Vec3,
        velocity: Vec3
    ) -> None:
        self.mass = mass
        self.position = position
        self.velocity = velocity

    def update_velocity(self,
        forces: list[ForceFunc],
        dt: float
    ) -> None:
        forces_values = [f(self) for f in forces]
        total_force = sum_vec3(forces_values)
        acceleration = total_force / self.mass
        self.velocity += acceleration * dt

    def update_position(self, dt: float) -> None:
        self.position += self.velocity * dt
\end{lstlisting}

\subsection{Forças}

Cada força é implementada como uma \emph{factory}: uma função que recebe os parâmetros físicos e retorna uma \texttt{ForceFunc}. Esse padrão - conhecido como \emph{closure} - permite que os parâmetros sejam fixados uma única vez na construção da força e capturados automaticamente em cada chamada dentro do laço de integração, sem a necessidade de passá-los explicitamente a cada passo.

\subsubsection{Força Gravitacional}

A força gravitacional atua exclusivamente na direção $-\hat{k}$ e tem magnitude proporcional à massa da partícula:
\begin{equation}
    \vec{F}_g(P) = -m\,a_g\,\hat{k}
\end{equation}

Apesar de simples, a força acessa $P.m$ a cada chamada, ilustrando que mesmo a gravidade depende do estado instantâneo - em particular da massa da partícula.

\begin{lstlisting}[caption={Força gravitacional como closure}, label={lst:grav}]
GRAVITATIONAL_ACCELERATION = 9.81   # [m/s^2]

def gravitational_force(
    a: float = GRAVITATIONAL_ACCELERATION
) -> ForceFunc:
    def f(p: Particle) -> Vec3:
        magnitude = -p.mass * a   # acessa p.mass do estado atual
        return Vec3(0, 0, magnitude)
    return f
\end{lstlisting}

\subsubsection{Força de Arrasto}

O arrasto possui coeficientes que decaem exponencialmente com a altitude e são assimétricos entre a subida e a descida. Os coeficientes de referência (subida) são:
\begin{equation}
    \beta_\uparrow(z) = \beta_0\,e^{-z/H}, \qquad
    \gamma_\uparrow(z) = \gamma_0\,e^{-z/H}
\end{equation}

e os coeficientes de descida são múltiplos desses:
\begin{equation}
    \beta_\downarrow(z) = r_\beta\,\beta_\uparrow(z), \qquad
    \gamma_\downarrow(z) = r_\gamma\,\gamma_\uparrow(z)
\end{equation}

A força resultante é \textbf{sempre oposta} ao vetor velocidade $\vec{v}$ nas três dimensões, com magnitude controlada pelos dois regimes:
\begin{equation}
    \vec{F}_d(P) = \begin{cases}
        -\bigl(\beta_\uparrow(z) + \gamma_\uparrow(z)\,\|\vec{v}\|\bigr)\,\vec{v}, & v_z > 0 \quad\text{(subida)}\\[6pt]
        -\bigl(\beta_\downarrow(z) + \gamma_\downarrow(z)\,\|\vec{v}\|\bigr)\,\vec{v}, & v_z \leq 0 \quad\text{(descida)}
    \end{cases}
\end{equation}

A condição $v_z > 0$ ou $v_z \leq 0$ define os \textbf{dois regimes dinâmicos da partícula}: subida e descida. Essa seleção é feita a cada passo de integração com base no estado instantâneo da partícula - não existe variável de fase explícita; a transição entre regimes é automática e determinada unicamente pelo sinal de $v_z$. Os parâmetros $r_\beta$ e $r_\gamma$ controlam a assimetria: quando $r_\beta > 1$, o arrasto linear na descida é mais intenso do que na subida, prolongando o tempo de queda em relação ao de ascensão.

A magnitude da força possui dois termos: o termo linear $\beta\,\|\vec{v}\|$ e o quadrático $\gamma\,\|\vec{v}\|^2$. Ambos dissipam energia e atuam na direção de $-\vec{v}$, sendo portanto válidos para as três componentes vetoriais simultaneamente.

\begin{lstlisting}[caption={Força de arrasto como closure - seleção de regime pelo estado $v_z$}, label={lst:ars}]
BETA_O  = 1.25  # [kg/s] Intensidade de arrasto linear
GAMMA_O = 0.1   # [kg/m] Intensidade de arrasto quadratico
R_BETA  = 5.0   # [adimensional] Fator de assimetria, arrasto linear
R_GAMMA = 2.5   # [adimensional] Fator de assimetria, arrasto quadratico
H = 2.0         # [m] Altura de escala

def drag_force(
    beta_o:  float = BETA_O,
    gamma_o: float = GAMMA_O,
    r_beta:  float = R_BETA,
    r_gamma: float = R_GAMMA,
    h:       float = H
) -> ForceFunc:
    def f(p: Particle) -> Vec3:
        speed    = p.velocity.length()            # ||v|| -- norma 3D
        beta_up  = beta_o  * exp(-p.position.z / h)
        gamma_up = gamma_o * exp(-p.position.z / h)

        if p.velocity.z > 0:   # subida
            beta  = beta_up
            gamma = gamma_up
        else:                   # descida ou apogeu
            beta  = r_beta  * beta_up
            gamma = r_gamma * gamma_up

        # sempre oposto a v, valido para as tres componentes
        return p.velocity * (-beta - gamma * speed)
    return f
\end{lstlisting}

\subsection{Integração Numérica}
\label{subsec:integracao}

A integração adota o método de Euler \emph{simplético}: a velocidade é atualizada \textbf{antes} da posição. Isso implica que $\vec{r}_{n+1}$ é calculado com a velocidade já corrigida $\vec{v}_{n+1}$, e não com $\vec{v}_n$ como no Euler explícito padrão. O livro \cite{kulp-pagonis-cm-2020} foi utilizado como primeira consulta de como implementar esse método de integração.

O estado $(t,\,\vec{r},\,\vec{v})$ é registrado \textbf{antes} de cada integração, garantindo que posição e velocidade armazenadas correspondam ao mesmo instante $t_n$ - consistência necessária para o cálculo correto das métricas. O laço termina quando $z < 0$, indicando que a partícula cruzou o solo.

\begin{lstlisting}[caption={Laço de integração - Euler simplético com salvamento de estado}, label={lst:sim}]
def simulate_until_ground(
    position: Vec3 | None = None,
    velocity: Vec3 | None = None,
    mass:     float = 1.0,
    a_g:      float = GRAVITATIONAL_ACCELERATION,
    beta_o:   float = BETA_O,
    gamma_o:  float = GAMMA_O,
    r_beta:   float = R_BETA,
    r_gamma:  float = R_GAMMA,
    h:        float = H,
    with_drag: bool = True
) -> SimulationResults:

    particle = Particle(mass=mass,
                        position=position or Vec3(0,0,0),
                        velocity=velocity or Vec3(0,0,30))

    forces = [gravitational_force(a=a_g)]
    if with_drag:
        forces.append(drag_force(beta_o, gamma_o, r_beta, r_gamma, h))

    t, i = 0.0, 0
    while particle.position.z >= 0 and i < MAX_STEPS:
        save_state(i, t)                          # salva estado ANTES de integrar
        particle.update_velocity(forces, TIME_STEP)  # atualiza v primeiro
        particle.update_position(TIME_STEP)          # usa v ja atualizado
        t += TIME_STEP
        i += 1

    return SimulationResults(
        times=times[:i],
        positions=positions[:i],
        velocities=velocities[:i],
        metrics=extract_metrics(times[:i], positions[:i], velocities[:i])
    )
\end{lstlisting}

\subsection{Extração de Métricas}

A partir dos arrays registrados, são extraídas quatro grandezas características da trajetória:

\begin{itemize}
    \item \textbf{Altura máxima} $z^* = \max_n z_n$: obtida pelo índice do máximo do array de posições verticais.
    \item \textbf{Tempo de subida} $\tau_\uparrow = t[\arg\max z]$: instante em que a partícula atinge $z^*$, marcando a transição entre os dois regimes dinâmicos.
    \item \textbf{Tempo de descida} $\tau_\downarrow = t_\text{final} - \tau_\uparrow$: diferença entre o último instante registrado (imediatamente antes de $z$ cruzar zero) e o instante do apogeu.
    \item \textbf{Velocidade de impacto} $\vec{v}^* = \vec{v}[N-1]$: último estado de velocidade salvo, que corresponde ao instante de maior proximidade ao solo antes da condição de parada ser ativada.
\end{itemize}

\begin{lstlisting}[caption={Extração de métricas a partir dos arrays simulados}, label={lst:metrics}]
def extract_metrics(times, positions, velocities) -> SimulationMetrics:
    z     = positions[:, 2]
    z_max = float(np.max(z))
    i_top = int(np.argmax(z))

    tau_up   = float(times[i_top])
    tau_down = float(times[-1]) - tau_up

    v_last  = velocities[-1]
    v_final = Vec3(float(v_last[0]),
                   float(v_last[1]),
                   float(v_last[2]))

    return SimulationMetrics(
        z_max=z_max,   tau_up=tau_up,
        tau_down=tau_down, v_final=v_final
    )
\end{lstlisting}
